\documentclass[UTF8,11pt]{ctexart}
\usepackage{amsmath,amssymb,geometry}
\geometry{margin=22mm}
\begin{document}
\section*{P8: 诊断样本八阶段运输的全灵活OST调度路径模型}
\paragraph{来源与范围}依据原论文PDF第12, 13, 14, 15, 16, 17, 18, 19, 20, 21, 22, 1, 26, 27页，(1)–(5)、(16)–(32)；以(2)–(3)闭合运输逻辑。
\subsection*{应用范围与参数}
选择论文的全灵活OST Scheduling and Routing MIP：在给定有限计划期$\mathcal T=\{1,\ldots,T\}$内联合优化每日路径与样本八阶段流转。

地点$V=\{1,\ldots,n+1\}$中，1为区级枢纽，$2,\ldots,n$为其他卫生机构，$L=n+1$为分子实验室；样本来源$I=\{1,\ldots,n\}$。有向弧$A=\{(i,j):i,j\in V,i\ne j\}$，对称行程距离$D_{ij}=D_{ji}\geq0$、$D_{ii}=0$。参数$c_t,H_t,R_t$分别为每日配送员数、每条路径的最大非枢纽停靠数和最大里程；全期/第一周里程预算为$B,B_1$，$\mathcal W_1=\{1,\ldots,\min(7,T)\}$。定义当日可服务地点$V_t=\{1\}\cup\{i\in V\setminus\{1\}:2D_{1i}\leq R_t\}$。所有涉及$V\setminus V_t$的弧和这些地点的访问量固定为零；顺序和距离辅助量仅为$V_t\setminus\{1\}$定义。来源集合$I$中的全部样本库存和流转均计入等待费用，包括当天不可访问来源的库存。$H_t\ge1$及$c_t\ge0$为整数。该日节点筛选是原式(26)服务半径必要条件在当前实现中的明确化。
为沿用原文(26)--(27)中以$R_t$为激活常数的距离约束，输入还须满足
$D_{ij}\leq D_{i1}+D_{1j}$（通过枢纽的三角不等式）。欧氏距离或同一无向网络上的最短路距离满足该条件。
原文第25页的合成实验明确使用欧氏距离；第33页说明运营距离矩阵曾由配送员记录修正。
该$R_t$形式适用于满足上述三角不等式的距离数据。


$u_i(t)$是已知每日样本到达；$\bar s_{ik}\geq0$为计划开始时各来源在阶段$k=1,\ldots,7$的存量；$\rho_k\geq0$为期末存量惩罚。阶段1为机构待取样本，2为枢纽待登记样本，3为枢纽已登记样本，4为实验室待检样本，5为实验室已有结果，6为枢纽待登记结果，7为枢纽已登记结果；阶段8表示已返回机构的完成样本，其等待费用为零。

对于处理步骤$k\in\{3,5,7\}$，$\alpha_k(\tau,t)\in[0,1]$给出第$\tau$日进入上一阶段的批次在第$t$日完成处理的比例，取$\alpha_k(\tau,t)=0$于$t<\tau$，允许$\alpha_k(t,t)>0$表示当天完成，同一批次总完成比例不超过1。$b_{ik}(t)\geq0$是起始处理中存量在$t$日的预定完成量，由计划前批次和相同处理日历确定，满足$\sum_t b_{ik}(t)\leq\bar s_{i,k-1}$；空初始系统时为0。该参数表示论文式(31)中$\tau\leq0$的历史流量及非空初始队列的预定完成量。

取样本数上界$M\geq\max_i[\sum_{k=1}^7\bar s_{ik}+\sum_tu_i(t)]$。各系数由实例给定。
\subsection*{变量}
$x_{ij}(t)$为选中弧，$y_i(t)$为地点访问指示；$h_i(t),g_i(t)$为非枢纽点的路径顺序和距离辅助量，$d(t)$为每日总里程。$s_{ik}(t)$为期末阶段存量，$f_{ik}(t)$为当日进入阶段$k=1,\ldots,8$的流量。流量连续非负，路径与访问变量为二元。
\subsection*{唯一完整模型}
\begin{equation}
\min\quad \frac1T\sum_{t=1}^T\sum_{i\in I}\sum_{k=1}^7s_{ik}(t)+\sum_{i\in I}\sum_{k=1}^7\rho_k s_{ik}(T).
\end{equation}
\begin{align}
d(t)&=\sum_{(i,j)\in A}D_{ij}x_{ij}(t) && t\in\mathcal T,\\
\sum_{t\in\mathcal T}d(t)&\leq B,&\sum_{t\in\mathcal W_1}d(t)&\leq B_1,\\
\sum_{j\ne1}x_{1j}(t)&=\sum_{j\ne1}x_{j1}(t)\leq c_t &&t\in\mathcal T,\\
\sum_{j\ne i}x_{ij}(t)&=\sum_{j\ne i}x_{ji}(t)=y_i(t) && i\ne1,t\in\mathcal T,\\
x_{1j}(t)&\leq y_1(t),\quad y_1(t)\leq\sum_{j\ne1}x_{1j}(t)&&j\ne1,t\in\mathcal T,\\
x_{jL}(t)&=x_{Lj}(t)=0&&j=2,\ldots,n,t\in\mathcal T,\\
0\leq h_i(t)&\leq H_t-1 && i\in V_t\setminus\{1\},t\in\mathcal T,\\
h_j(t)-h_i(t)&\geq1-H_t(1-x_{ij}(t))&& i,j\in V_t\setminus\{1\},i\ne j,t\in\mathcal T,\\
D_{i1}\leq g_i(t)&\leq R_t-D_{1i}&& i\in V_t\setminus\{1\},t\in\mathcal T,\\
g_j(t)-g_i(t)&\geq D_{ij}-R_t(1-x_{ij}(t))&&i,j\in V_t\setminus\{1\},i\ne j,t\in\mathcal T.
\end{align}
访问指示按原文第13页的定义$y_i=\max_jx_{ij}$闭合，非枢纽一天至多访问一次按第18页文字明确列出；枢纽来源$i=1$使用相同访问定义，其样本计入目标和流转。

设触发访问量$z_{ik}(t)=y_i(t)$于$k\in\{2,8\}$，$z_{ik}(t)=y_L(t)$于$k\in\{4,6\}$，它仅是表达式缩写。对所有$i\in I,t\in\mathcal T,k\in\{2,4,6,8\}$，写出$f_{ik}(t)=z_{ik}(t)s_{i,k-1}(t-1)$的完整线性化：
\begin{align}
0\leq f_{ik}(t)&\leq M z_{ik}(t),\\
f_{ik}(t)&\leq s_{i,k-1}(t-1),\\
f_{ik}(t)&\geq s_{i,k-1}(t-1)-M(1-z_{ik}(t)).
\end{align}
\begin{align}
s_{ik}(0)&=\bar s_{ik} && i\in I,k=1,\ldots,7,\\
f_{i1}(t)&=u_i(t) &&i\in I,t\in\mathcal T,\\
f_{ik}(t)&=b_{ik}(t)+\sum_{\tau=1}^{t}\alpha_k(\tau,t)f_{i,k-1}(\tau)&&i\in I,k\in\{3,5,7\},t\in\mathcal T,\\
s_{ik}(t)&=s_{ik}(t-1)+f_{ik}(t)-f_{i,k+1}(t)&&i\in I,k=1,\ldots,7,t\in\mathcal T,\\
s_{ik}(t),f_{ik}(t),d(t)&\geq0,\quad x_{ij}(t),y_i(t)\in\{0,1\}&&\text{所有已定义下标}.
\end{align}
原文(29)--(30)印出的两侧不等式未单独保证未访问时零流量；本文件依据明确的(2)--(3)乘积定义补齐逻辑。其余可选访问频率政策未默认启用：它们是特定运营方的附加限制，非基础(32)模型必要条件。此模型是有限期确定性计划，不承诺其单次最优解是无限期随机控制最优策略。

\end{document}
